% !TeX root = ../main.tex
\chapter{MỞ ĐẦU}
\label{chap:modau}

\section{Bối cảnh nghiên cứu}
\label{sec:boicanh}

Sầu riêng (\textit{Durio zibethinus}) là một trong những cây ăn quả nhiệt đới có
giá trị kinh tế cao nhất khu vực Đông Nam Á. Tại Việt Nam, cây sầu riêng được
canh tác tập trung ở Tây Nguyên, Đông Nam Bộ và đồng bằng sông Cửu Long, đóng vai
trò quan trọng trong cơ cấu thu nhập của nhiều hộ nông dân. Đặc điểm của loại cây
này là chu kỳ kiến thiết cơ bản dài --- thường phải mất từ bốn đến sáu năm kể từ
khi trồng mới cho thu hoạch ổn định --- nên mỗi tổn thất do sâu bệnh gây ra đều
kéo theo hệ quả kinh tế lâu dài, không thể khắc phục trong một vụ.

Trong các nhóm bệnh hại trên cây sầu riêng, bệnh trên lá chiếm vị trí đặc biệt vì
hai lý do. Thứ nhất, lá là cơ quan quang hợp; diện tích lá bị tổn thương giảm
trực tiếp làm suy giảm khả năng tích luỹ chất khô, ảnh hưởng đến năng suất và
chất lượng quả. Thứ hai, triệu chứng trên lá thường xuất hiện \emph{sớm hơn}
triệu chứng trên thân và quả, nên lá đóng vai trò chỉ báo cho phép can thiệp kịp
thời. Việc phát hiện chính xác loại bệnh và xác định được phạm vi tổn thương ngay
từ giai đoạn đầu vì thế có ý nghĩa quyết định đối với hiệu quả phòng trừ.

\subsection{Nguồn gốc bộ dữ liệu}
\label{subsec:nguongoc}

Bộ dữ liệu được sử dụng trong luận văn gồm \num{4437} ảnh lá sầu riêng thuộc năm
lớp: bốn lớp bệnh và một lớp lá khoẻ. Chi tiết về phân bố lớp, đặc trưng màu sắc
và cách phân chia tập dữ liệu được trình bày trong \cref{sec:dulieu}.

Ảnh được thu thập trực tiếp tại các vườn trồng sầu riêng ở khu vực Đông Nam Á.
Nhãn ở mức ảnh --- tức việc xác định mỗi ảnh thuộc lớp bệnh nào --- do cán bộ
chuyên môn nông nghiệp thực hiện và thẩm định; nhóm nghiên cứu không tự gán nhãn
dựa trên quan sát của người không có chuyên môn bệnh học thực vật. Đây là cơ sở
để xem nhãn mức ảnh của bộ dữ liệu là đáng tin cậy, và cũng là tiền đề cho toàn
bộ quy trình học giám sát yếu được đề xuất: nếu nhãn mức ảnh sai thì bản đồ chú ý
sinh ra từ nó, và do đó cả nhãn giả lẫn mô hình phân đoạn, đều mất giá trị.

Bộ dữ liệu và trọng số mô hình được chia sẻ theo yêu cầu hợp lý gửi tới tác giả
liên hệ; luận văn không công bố công khai dữ liệu thô.

Cần lưu ý ba đặc điểm về trạng thái của dữ liệu khi được đưa vào nghiên cứu này,
vì chúng ảnh hưởng trực tiếp đến phạm vi kết luận có thể rút ra:

\begin{itemize}
    \item \textbf{Ảnh đã được tiền xử lý trước.} Toàn bộ \num{4437} ảnh đều ở
          đúng kích thước $224 \times 224$ điểm ảnh với tỉ lệ khung hình bằng 1
          và độ lệch chuẩn bằng 0 --- nghĩa là việc cắt và co giãn đã được thực
          hiện ở khâu trước. Hệ quả là luận văn không thể khảo sát ảnh hưởng của
          độ phân giải đầu vào lên chất lượng phân đoạn.
    \item \textbf{Siêu dữ liệu EXIF đã bị loại bỏ.} Kiểm tra trực tiếp các tệp
          ảnh cho thấy không còn thông tin về thời điểm chụp, thiết bị chụp hay
          toạ độ định vị. Vì vậy không thể phân tích ảnh hưởng của điều kiện chụp
          --- ánh sáng, thiết bị, mùa vụ --- lên hiệu năng mô hình.
    \item \textbf{Cấu trúc tên tệp phản ánh quy trình gán nhãn.} Các tệp được đặt
          tên theo dạng \texttt{to\_\allowbreak label\_\allowbreak N.\allowbreak jpg} với chỉ số $N$ đánh liên tục
          xuyên suốt các lớp, cho thấy dữ liệu ban đầu tồn tại dưới dạng một kho
          ảnh chung chờ gán nhãn, sau đó mới được phân vào các thư mục theo lớp.
\end{itemize}

\TODO{Bổ sung nếu lấy được từ nhóm thu thập dữ liệu gốc: (1) địa phương cụ thể
và thời gian thu thập; (2) thiết bị chụp và độ phân giải ảnh gốc trước khi co về
$224\times224$; (3) số lượng chuyên gia tham gia gán nhãn và mức độ đồng thuận
giữa họ nếu có nhiều hơn một người; (4) giấy phép sử dụng dữ liệu. Bốn thông tin
này không tồn tại trong kho mã nguồn và cũng không được ghi trong công bố liên
quan của nhóm.}

Phương pháp chẩn đoán truyền thống dựa hoàn toàn vào quan sát trực tiếp của cán
bộ kỹ thuật nông nghiệp. Cách làm này bộc lộ ba hạn chế cố hữu: tốn nhiều công
lao động khi diện tích canh tác lớn; kết quả mang tính chủ quan và thiếu nhất
quán giữa những người quan sát khác nhau; và đặc biệt khó phân biệt các bệnh có
biểu hiện hình thái tương tự nhau. Những hạn chế này tạo ra nhu cầu thực tế về
một công cụ chẩn đoán tự động, khách quan và có khả năng triển khai ở quy mô lớn.

Học sâu (\en{deep learning}) đã chứng minh hiệu quả trong bài toán nhận dạng bệnh
cây trồng qua nhiều nghiên cứu, từ công trình nền tảng của Mohanty và cộng
sự~\cite{mohanty2016plantdisease} trên bộ dữ liệu PlantVillage, đến các khảo sát
tổng quan của Kamilaris và Prenafeta-Boldú~\cite{kamilaris2018survey} hay
Ferentinos~\cite{ferentinos2018plantdisease}. Tuy nhiên, phần lớn các nghiên cứu
này dừng lại ở mức \emph{phân loại ảnh}: mô hình trả lời câu hỏi ``ảnh này thuộc
bệnh nào'' nhưng không trả lời được câu hỏi ``vùng bệnh nằm ở đâu và rộng bao
nhiêu''. Đây chính là khoảng trống mà luận văn hướng tới.

\section{Phát biểu vấn đề}
\label{sec:vande}

Luận văn giải quyết ba vấn đề gắn bó chặt chẽ với nhau.

\subsection{Chi phí gán nhãn ở mức điểm ảnh}

Để huấn luyện một mô hình phân đoạn (\en{segmentation}) theo cách giám sát đầy đủ,
mỗi ảnh cần một mặt nạ nhị phân (\en{binary mask}) chỉ rõ từng điểm ảnh thuộc vùng
bệnh hay vùng lành. Công việc khoanh vùng thủ công này đòi hỏi chuyên gia bệnh
học thực vật, tốn khoảng vài phút cho mỗi ảnh, và với bộ dữ liệu \num{4437} ảnh
thì tổng chi phí trở nên khó chấp nhận trong điều kiện nghiên cứu thông thường.
Trong khi đó, nhãn ở \emph{mức ảnh} --- tức chỉ ghi nhãn ``ảnh này bị bệnh
Phomopsis'' --- lại rẻ hơn nhiều bậc và thường đã có sẵn.

Bài toán đặt ra là: \emph{liệu có thể tận dụng nhãn mức ảnh vốn rẻ để tạo ra
nhãn mức điểm ảnh mà không cần khoanh vùng thủ công hay không?} Đây là phạm trù
của học giám sát yếu (\en{weakly-supervised learning}).

\subsection{Tính khả diễn giải của mô hình}

Mạng nơ-ron tích chập vận hành như một hộp đen: người dùng nhận được nhãn dự đoán
kèm độ tin cậy nhưng không biết mô hình căn cứ vào đặc trưng nào. Trong bối cảnh
nông nghiệp, nơi một khuyến cáo phun thuốc sai có thể gây thiệt hại trực tiếp,
sự thiếu minh bạch này là rào cản thực sự đối với việc chấp nhận công nghệ. Cán
bộ kỹ thuật cần nhìn thấy mô hình ``đang nhìn vào đâu'' để đối chiếu với kinh
nghiệm chuyên môn trước khi tin vào kết luận của nó.

\subsection{Nhu cầu định vị chính xác vùng tổn thương}

Nhãn phân loại cho biết \emph{có} bệnh, nhưng quyết định canh tác lại cần biết
\emph{mức độ} bệnh. Tỉ lệ diện tích lá bị tổn thương là căn cứ để lựa chọn giữa
theo dõi tiếp, phun thuốc cục bộ hay cắt tỉa. Chỉ mô hình phân đoạn mới cung cấp
được đại lượng này.

Ba vấn đề trên gợi ý một hướng tiếp cận thống nhất: nếu bản đồ chú ý của mô hình
phân loại có thể được trực quan hoá để giải quyết vấn đề thứ hai, thì chính bản
đồ đó cũng có thể được chuyển hoá thành mặt nạ phân đoạn để giải quyết vấn đề thứ
nhất và thứ ba. Đây là ý tưởng trung tâm của luận văn.

\section{Mục tiêu và câu hỏi nghiên cứu}
\label{sec:muctieu}

\subsection{Mục tiêu tổng quát}

Xây dựng và đánh giá một quy trình học giám sát yếu hoàn chỉnh cho bài toán phát
hiện bệnh lá sầu riêng, đi từ phân loại ảnh, qua giải thích bằng bản đồ kích
hoạt, đến sinh nhãn giả và huấn luyện mô hình phân đoạn --- toàn bộ mà không sử
dụng bất kỳ nhãn thủ công nào ở mức điểm ảnh.

\subsection{Mục tiêu cụ thể}

\begin{enumerate}
    \item Huấn luyện mô hình phân loại \EffNet{} cho năm lớp bệnh lá sầu riêng và
          đánh giá đầy đủ trên tập kiểm thử độc lập.
    \item Áp dụng \GradCAM{} và \GradCAMpp{} đa tỉ lệ để trực quan hoá vùng chú ý
          của mô hình phân loại, đồng thời kiểm chứng rằng mô hình học đặc trưng
          bệnh học có ý nghĩa chứ không phải học nhiễu nền.
    \item Thiết kế quy trình chuyển hoá bản đồ nhiệt thành mặt nạ nhị phân, với
          cơ chế kiểm soát độ phủ theo từng lớp bệnh.
    \item Khảo sát mô hình nền tảng \SAM{} trong vai trò bộ tinh chỉnh biên
          (\en{label refiner}) cho nhãn giả.
    \item Huấn luyện và so sánh ba phiên bản mô hình phân đoạn tương ứng với ba
          thế hệ nhãn giả, phân tích tách bạch ảnh hưởng của từng yếu tố.
    \item Chỉ rõ giới hạn phương pháp luận của việc đánh giá phân đoạn khi không
          có nhãn chuẩn (\en{ground truth}) ở mức điểm ảnh.
\end{enumerate}

\subsection{Câu hỏi nghiên cứu}

\begin{itemize}
    \item[\textbf{CH1.}] Mô hình phân loại \EffNet{} đạt độ chính xác bao nhiêu
          trên năm lớp bệnh lá sầu riêng, và những cặp lớp nào dễ bị nhầm lẫn
          nhất?
    \item[\textbf{CH2.}] Bản đồ chú ý sinh ra từ mô hình phân loại có tập trung
          vào vùng tổn thương thực sự hay không, và mức độ tin cậy của chúng khác
          nhau thế nào giữa các lớp bệnh?
    \item[\textbf{CH3.}] Việc tinh chỉnh biên bằng \SAM{} có cải thiện chất lượng
          nhãn giả, và qua đó cải thiện mô hình phân đoạn, hay không?
    \item[\textbf{CH4.}] Trong điều kiện không có nhãn chuẩn ở mức điểm ảnh, các
          chỉ số IoU và Dice giữa những phiên bản khác nhau của quy trình có so
          sánh trực tiếp được với nhau không? Nếu không thì phép so sánh nào là
          hợp lệ?
\end{itemize}

Câu hỏi CH4 là câu hỏi mang tính phương pháp luận và, như sẽ thấy ở
\cref{chap:thucnghiem}, câu trả lời cho nó làm thay đổi hoàn toàn cách diễn giải
các con số của ba câu hỏi còn lại.

\section{Đối tượng, phạm vi và giới hạn}
\label{sec:phamvi}

\subsection{Đối tượng nghiên cứu}

Đối tượng là ảnh lá sầu riêng đơn lẻ ở độ phân giải $224 \times 224$ điểm ảnh,
thuộc năm lớp: Algal Leaf Spot (đốm tảo), Allocaridara Attack (rầy phấn tấn
công), Healthy Leaf (lá khoẻ), Leaf Blight (cháy lá) và Phomopsis Leaf Spot (đốm
lá Phomopsis).

\subsection{Phạm vi}

\begin{itemize}
    \item \textbf{Bài toán phân loại}: đa lớp, năm lớp, một nhãn trên mỗi ảnh.
    \item \textbf{Bài toán phân đoạn}: nhị phân --- phân biệt vùng bệnh với phần
          còn lại. Luận văn \emph{không} thực hiện phân đoạn đa lớp, tức không
          phân biệt loại bệnh ở mức điểm ảnh.
    \item \textbf{Dữ liệu}: chỉ sử dụng bộ dữ liệu \num{4437} ảnh nêu trên; không
          thu thập bổ sung và không sử dụng dữ liệu ngoài.
    \item \textbf{Mô hình}: các kiến trúc có kích thước vừa phải, huấn luyện được
          trên một GPU đơn.
\end{itemize}

\subsection{Giới hạn cần lưu ý trước khi đọc kết quả}
\label{subsec:gioihan-quantrong}

Có một giới hạn mang tính bản chất cần được nêu ngay từ đầu, vì nó chi phối cách
diễn giải toàn bộ kết quả phân đoạn trong luận văn:

\begin{quote}
\textbf{Bộ dữ liệu không có nhãn chuẩn ở mức điểm ảnh.} Vì vậy, mọi chỉ số IoU và
Dice được báo cáo trong \cref{chap:thucnghiem} đo mức độ mô hình phân đoạn
\emph{tái tạo lại được bộ nhãn giả} mà chính nó được huấn luyện trên đó --- chứ
\emph{không} đo độ chính xác so với vùng bệnh thật.
\end{quote}

Hệ quả trực tiếp của giới hạn này là: các phiên bản sử dụng bộ nhãn giả khác nhau
sẽ được đánh giá trên những mốc tham chiếu khác nhau, nên giá trị IoU của chúng
\emph{không so sánh trực tiếp được với nhau}. Luận văn không né tránh mà phân
tích thẳng vào vấn đề này ở \cref{sec:caveat}, đồng thời chỉ rõ phép so sánh nào
vẫn giữ được tính hợp lệ nội tại.

Ngoài ra, luận văn có ba giới hạn thứ cấp: (i) các hạn chế bắt nguồn từ trạng
thái dữ liệu đầu vào đã nêu ở \cref{subsec:nguongoc} --- ảnh đã được co về
$224 \times 224$ và siêu dữ liệu EXIF đã bị loại bỏ --- khiến không thể khảo sát
ảnh hưởng của độ phân giải cũng như điều kiện chụp; (ii) các ngưỡng độ phủ theo
lớp được chọn theo suy luận bệnh học chứ chưa được hiệu chỉnh dựa trên nhãn
chuẩn; (iii) hệ thống được đánh giá trên ảnh chụp lá đơn lẻ, chưa kiểm chứng
trên ảnh chụp toàn tán trong điều kiện đồng ruộng.

\section{Đóng góp của luận văn}
\label{sec:donggop}

\begin{enumerate}
    \item \textbf{Một quy trình học giám sát yếu bốn giai đoạn hoàn chỉnh} cho
          bệnh lá sầu riêng: phân loại $\rightarrow$ giải thích bằng \GradCAMpp{}
          $\rightarrow$ tinh chỉnh biên bằng \SAM{} $\rightarrow$ phân đoạn bằng
          \UNetpp{}, được cài đặt đầy đủ và có thể tái tạo.

    \item \textbf{Cơ chế ngưỡng phân vị theo lớp} (\en{per-class percentile
          thresholding}) cho phép kiểm soát trực tiếp độ phủ của mặt nạ theo đặc
          điểm từng bệnh, thay cho ngưỡng Otsu vốn giả định phân bố hai đỉnh và
          gây phân đoạn thừa khi bản đồ nhiệt được hợp nhất đa tỉ lệ. Cơ chế này
          đạt độ phủ mục tiêu với sai lệch dưới 0,4\% cho mọi lớp.

    \item \textbf{Phân tích định lượng về \SAM{} trong vai trò bộ tinh chỉnh
          nhãn}, bao gồm việc chỉ ra và giải thích một chế độ thất bại chưa được
          nhấn mạnh trong tài liệu: cơ chế bảo vệ dựa trên IoU chỉ chặn được
          trường hợp \SAM{} \emph{trôi} sang vùng khác, mà không chặn được trường
          hợp \SAM{} \emph{nở} mặt nạ ra bao trọn cả chiếc lá. Hiện tượng nở mặt
          nạ được đo cụ thể: độ phủ trung bình tăng từ 15,04\% lên 25,43\%.

    \item \textbf{Một phân tích tách bạch (\en{ablation}) trung thực} về khả năng
          so sánh của các chỉ số phân đoạn khi thiếu nhãn chuẩn. Luận văn chỉ ra
          rằng phiên bản có chỉ số cao nhất về mặt số học lại không phải phiên bản
          định vị bệnh tốt nhất, và giải thích tại sao --- một kết luận có giá trị
          cảnh báo cho các nghiên cứu học giám sát yếu tương tự.

    \item \textbf{Bộ tài liệu kiểm kê số liệu} (\texttt{INVENTORY.md}) truy vết
          mọi con số trong luận văn về đúng notebook và ô lệnh đã sinh ra nó, kèm
          danh sách các điểm chưa đủ bằng chứng. Cách làm này bảo đảm tính minh
          bạch và tái lập của kết quả.
\end{enumerate}

\section{Bố cục luận văn}
\label{sec:bocuc}

Phần còn lại của luận văn được tổ chức như sau.

\textbf{\Cref{chap:cosolythuyet}} trình bày cơ sở lý thuyết: mạng nơ-ron tích
chập và học chuyển giao, kiến trúc U-Net cùng biến thể \UNetpp{}, các hàm mất mát
cho dữ liệu mất cân bằng, họ phương pháp \GradCAM{}, khung học giám sát yếu, mô
hình nền tảng \SAM{}, và các kỹ thuật xử lý ảnh cổ điển được sử dụng.

\textbf{\Cref{chap:phuongphap}} mô tả phương pháp đề xuất: tổng quan quy trình
bốn giai đoạn, phân tích khám phá bộ dữ liệu, thiết kế mô hình phân loại, và ba
thế hệ nhãn giả V1--V3 cùng lý do kỹ thuật của từng cải tiến.

\textbf{\Cref{chap:thucnghiem}} trình bày kết quả thực nghiệm: hiệu năng phân
loại, phân tích định tính bản đồ chú ý, chất lượng nhãn giả qua các phiên bản,
hiện tượng nở mặt nạ của \SAM{}, kết quả phân đoạn, phân tích tách bạch và phân
tích lỗi.

\textbf{\Cref{chap:ketluan}} tổng kết đóng góp, trả lời trực tiếp bốn câu hỏi
nghiên cứu, thảo luận hạn chế và đề xuất hướng phát triển.

Bốn phụ lục cung cấp bảng siêu tham số đầy đủ, nhật ký huấn luyện chi tiết, mô tả
ứng dụng minh hoạ và hướng dẫn tái tạo toàn bộ kết quả.
